\section{Related Work}\label{sec:related}

\subsection{Formalized Complexity and Mechanized Proof Infrastructure}

Formalizing complexity-theoretic arguments in proof assistants remains comparatively sparse. Lean and Mathlib provide the current host environment for this artifact \cite{Lean2015,demoura2021lean4,mathlib2020}. Closest mechanized precedents include verified computability/complexity developments in Coq and Isabelle \cite{forster2019verified, kunze2019formal, nipkow2002isabelle, nipkow2014concrete, haslbeck2016verified}, and certifying toolchain work that treats proofs as portable machine-checkable evidence \cite{leroy2009compcert, necula1997proof}.

Recent LLM-evaluation work also highlights inference nondeterminism from numerical precision and hardware configuration, reinforcing the value of proof-level claims whose validity does not depend on stochastic inference trajectories \cite{yuan2025fp32}.

\subsection{Computational Decision Theory}

The complexity of decision-making has been studied extensively. Foundational treatments of computational complexity and strategic decision settings establish the baseline used here \cite{papadimitriou1994complexity,arora2009computational}. Our work extends this to the meta-question of identifying relevant information. Decision-theoretic and information-selection framing in this paper also sits in the tradition of statistical sufficiency and signaling/information economics \cite{fisher1922mathematical, spence1973job, myerson1979incentive, milgrom1986price, kamenica2011bayesian}.

\paragraph{Categorical framing of the quotient object.}
The universal property of the optimizer quotient object is not, by itself, a probabilistic construction. In the category \textbf{Set}, quotienting states by equality of $\Opt$ is the standard coimage construction for the optimizer map $\Opt : S \to \mathcal{P}(A)$, canonically equivalent to its image/range \cite{maclane1998categories}. The novelty in this paper is therefore not that coimages exist, but that this particular coimage organizes coordinate sufficiency, witness lower bounds, and physical collapse constraints in one mechanized chain.\claimstamp{T\ref{thm:quotient-universal}}{DM}\leanmeta{\LH{QT1}, \LH{QT3}, \LH{QT5}.}

Large-sample Bayesian-network structure learning and causal identification are already known to be hard in adjacent formulations \cite{chickering2004large,shpitser2006identification,koller2009probabilistic}. Our object differs: we classify coordinate sufficiency for optimal-action invariance across static/stochastic/sequential regimes, with theorem-level placements and mechanized reductions.

\paragraph{Relation to prior hardness results.}
Closest adjacent results are feature-selection/model-selection hardness results for predictive subset selection~\cite{blum1997selection, amaldi1998complexity}. Relative to those works, this paper changes two formal objects: (i) the reductions are machine-checked (TAUTOLOGY and $\exists\forall$-SAT mappings with explicit polynomial bounds), and (ii) the output is a regime-typed hardness/tractability map for decision relevance under explicit encoding assumptions. The target predicate here is decision relevance, not predictive compression.

\paragraph{Functional information and selection.}
Prior work defines function-relative information objects in physical systems. Szostak defines functional information as the negative log of the fraction of configurations achieving a specified function~\cite{szostak2003functional}. Wong et al.\ study function and selection in evolving systems and state an increasing-functional-information law under selection~\cite{wong2023roles}. This paper addresses a complementary object: the computational decision predicate for when projected coordinates are sufficient to preserve optimal-action correspondence, with explicit complexity placements and mechanized proofs.

\subsection{Succinct Representations and Regime Separation}

Representation-sensitive complexity is established in planning and decision-process theory: classical and compactly represented MDP/planning problems exhibit sharp complexity shifts under different input models~\cite{papadimitriou1987mdp,mundhenk2000mdp,littman1998probplanning}. The explicit-vs-succinct separation in this paper is the same representation-sensitive phenomenon for the coordinate-sufficiency predicate.

The distinction in this paper is the object and scope of the classification: we classify \emph{decision relevance} (sufficiency, anchor sufficiency, and minimum sufficient sets) for a fixed decision relation, with theorem-level regime typing and mechanized reduction chains.

\subsection{Oracle and Query-Access Lower Bounds}

Query-access lower bounds are a standard source of computational hardness transfer~\cite{dobzinski2012query,buhrman2002complexity}. Our \texttt{[Q\_fin]} and \texttt{[Q\_bool]} results are in this tradition, but specialized to the same sufficiency predicate used throughout the paper: they establish access-obstruction lower bounds while keeping the underlying decision relation fixed across regimes.

Complementary companion work studies zero-error identification under constrained observation and proves rate-query lower bounds plus matroid structure of minimal distinguishing query sets~\cite{simas2026identification}. Our object here is different: sufficiency of coordinates for optimal-action invariance in a decision problem, with regime-map placement (coNP/PP/PSPACE), static inner decomposition (coNP/$\Sigma_2^P$ decision layer plus access/structure layers), and regime-typed transfer theorems.

\subsection{Energy Landscapes and Chemical Reaction Networks}

The geometric interpretation of decision complexity connects to the physics of energy landscapes. Wales established a widely used picture in which molecular configurations form high-dimensional landscapes with basins and transition saddles~\cite{wales2003energy,wales2005energy,wales2018exploring}. Protein-folding theory uses the same landscape language via funnel structure toward native states~\cite{onuchic1997protein}. In this paper's formal model, the corresponding core object is coordinate-projection invariance of the optimizer map: sufficiency asks whether fixing coordinates preserves \(\Opt\), while insufficiency is witnessed by agreeing projections with different optimal sets (Definition~\ref{def:sufficient}, Proposition~\ref{prop:insufficiency-counterexample}).\claimstamp{D\ref{def:sufficient};P\ref{prop:insufficiency-counterexample}}{DM}\leanmeta{\LHrng{DP}{7}{8}.}

Section~\ref{sec:physical-transport} makes this transfer rule theorem-level: encoded physical claims are derived from core claims through an explicit assumption-transfer map, and an encoded physical counterexample refutes the corresponding core rule on that encoded slice.\claimstamp{P\ref{prop:physical-claim-transport};C\ref{cor:physical-counterexample-core-failure}}{AR}\leanmeta{\LHrng{CT}{5}{8}.}

Under that mapping, the explicit-vs-succinct split is the mechanized statement of access asymmetry: [E] admits polynomial verification on explicit state access, while [S+ETH] admits worst-case exponential lower bounds on succinct inputs (Theorem~\ref{thm:dichotomy}); finite query access already exhibits indistinguishability barriers (Proposition~\ref{prop:query-regime-obstruction}). The complexity separation is mechanized.\claimstamp{T\ref{thm:dichotomy};P\ref{prop:query-regime-obstruction}}{E,S+ETH,Q\_fin}\leanmeta{\LH{CC16}, \LHrng{CC}{23}{24}, \LH{CC50}.}

The geometric core is mechanized directly: product-space slice cardinality, hypercube node-count identity, and node-vs-edge separation are proved in Lean at the decision abstraction level, with edge witnesses tied exactly to non-sufficiency.\leanmeta{\LHrng{DG}{5}{6}, \LHrng{DG}{15}{16}, \LH{DG18}.}

The circuit-to-decision instantiation is also mechanized in a separate bridge layer: geometry and dynamics are represented as typed components of a circuit object; decision circuits add an explicit interpretation layer; molecule-level constructors instantiate both views.\leanmeta{\LHrng{IN}{4}{6}, \LHrng{IN}{13}{16}.}

Chemical reaction networks provide a concrete physical class where this encoding applies directly. Prior work establishes hardness for output optimization, reconfiguration, and reconstruction in CRN settings~\cite{andersen2012maximizing,alaniz2023complexity,flamml2013complexity}. Those domain-specific CRN complexity results are external literature claims, not mechanized in this Lean artifact. The mechanized statement here is the transfer principle: once a CRN decision task is encoded as a C1--C4 decision problem, it inherits the same sufficiency hardness/tractability regime map (Theorems~\ref{thm:config-reduction}, \ref{thm:sufficiency-conp}, \ref{thm:tractable}).\claimstamp{T\ref{thm:config-reduction},\ref{thm:sufficiency-conp},\ref{thm:tractable}}{E,S,TA}\leanmeta{\LH{CR1}, \LH{CC61}, \LHrng{CC}{69}{72}.}

The physical bridge is not a single-domain add-on: Theorem~\ref{thm:physical-bridge-bundle} composes law-objective semantics, reduction equivalence, and hardness lower-bound accounting in one mechanized theorem, so physical instantiation and complexity claims share one proof chain.\claimstamp{T\ref{thm:physical-bridge-bundle}}{AR}\leanmeta{\LH{CT4}.}

For matter-only dynamics, the law-induced objective specialization is mechanized for arbitrary universes (arbitrary state/action types and transition-feasibility relations): the utility schema is induced from transition laws, and under a strict allowed-vs-blocked gap with nonempty feasible set, the optimizer equals the feasible-action set; with logical determinism at the action layer, the optimizer is singleton.\claimstamp{P\ref{prop:law-instance-objective-bridge}}{AR}\leanmeta{\LHrng{DQ}{5}{8}.}

The interface-time semantics are also explicit and mechanized: timed states are $\mathbb{N}$-indexed, one decision event is exactly one tick, run length equals elapsed interface time, and this unit-time law is invariant across substrate tags under interface-consistency assumptions.\claimstamp{P\ref{prop:time-discrete},\ref{prop:decision-unit-time},\ref{prop:run-time-accounting},\ref{prop:substrate-unit-time}}{AR}\leanmeta{\LH{DT6}, \LH{DT8}, \LH{DT10}, \LH{DT13}, \LH{DT16}, \LH{DT24}.}

Regime-side interface machinery is also formalized as typed access objects with an explicit simulation preorder and certificate-upgrade statements. The access-channel and five-way composition theorems are encoded as assumption-explicit composition laws.\claimstamp{P\ref{prop:physical-claim-transport};T\ref{thm:physical-bridge-bundle}}{AR}\leanmeta{\LHrng{ARG}{2}{3}, \LH{ARG6}, \LHrng{ARG}{12}{13}, \LH{ARG17}, \LH{ARG19}.}

\subsection{Thermodynamic Costs and Landauer's Principle}

The thermodynamic lift in Section~\ref{sec:dichotomy} converts bit-operation lower bounds into energy/carbon lower bounds via linear conversion constants. In this artifact, those conversion consequences are mechanized with explicit premises through Propositions~\ref{prop:thermo-lift}, \ref{prop:thermo-hardness-bundle}, \ref{prop:thermo-mandatory-cost}, and \ref{prop:thermo-conservation-additive}. Landauer's principle and reversible-computation context~\cite{landauer1961irreversibility,bennett2003notes} define a minimum floor for irreversible bit erasure, not a generally tight expression for the actual energetic cost of constrained computation. The core stochastic-thermodynamics review literature~\cite{vandenbroeck2015ensemble} supplies the ensemble- and trajectory-level entropy-production framework used to state such refinements. Wolpert et al.\ \cite{wolpert2024stochastic} identify the relevant constrained-computation regimes: periodic driving, modular organization, finite-time operation, and input mismatch. Manzano et al.\ \cite{manzano2024absolute} make stopping-time and absolute-irreversibility effects explicit and derive universal dissipation bounds in that setting. Yadav, Yousef, and Wolpert \cite{yadav2026minimal} connect mismatch cost to repeated circuit operation and structural resource lower bounds. Our formal model now represents those inputs at three levels. The coarse interface \LH{WC1}--\LH{WC5} records the Landauer floor plus a generic nonnegative overhead. The refined decomposition \LH{WP1}--\LH{WP9} separates that overhead into mismatch and residual terms. Within that refined layer, the mismatch branch is no longer only imported: \LH{WM1}--\LH{WM4} derive nonnegativity, exact vanishing, strict positivity under explicit finite-distribution mismatch, and the conservative nat-valued lower-bound bridge from the existing KL machinery, while \LH{WM5}--\LH{WM6} inject that derived branch into the Wolpert decomposition and prove strict separation above the Landauer floor. The finite discrete residual branch is now also theorem-level and exhaustive at the local edge level: \LH{WR1}--\LH{WR3} give the bidirectional asymmetry witness, \LH{WR6} performs the reverse-edge case split, \LH{WR7}--\LH{WR8} inject the resulting unified finite residual witness into the Wolpert decomposition, and \LH{WR9}--\LH{WR10} repackage the same argument as a single process-level witness theorem for finite computational-state systems. Broader stopping-time / absolute-irreversibility residual regimes remain represented by the separate cited premise \LH{WP5}. This keeps the causal structure of the cited work explicit while limiting the mechanization claim to what this artifact actually proves: the mismatch branch is derived, a finite discrete residual branch is derived by an exhaustive edge split and repackaged at process level, broader stopping-time residual regimes remain imported, and the composition layer is fully machine-checked.\claimstamp{P\ref{prop:thermo-lift},\ref{prop:thermo-hardness-bundle},\ref{prop:thermo-mandatory-cost},\ref{prop:thermo-conservation-additive}}{AR}\leanmeta{\LHrng{CC}{64}{68}, \LHrng{WC}{1}{5}, \LHrng{WM}{1}{6}, \LHrng{WR}{1}{10}, \LHrng{WP}{1}{3}, \LHrng{WP}{5}{9}.}

\paragraph{Physics assumptions decompose into regime tags.}
The paper does not create domain-specific tags for physics assumptions. Instead, every physics assumption decomposes into the same regime tags used throughout the paper. This uniformity states that physics assumptions are not epistemically privileged; they are declared assumptions, tracked with the same discipline as the paper's complexity-theoretic assumptions.
The explicit physical-to-core mapping does not make the results less physical. It standardizes premise exposure and transfer obligations: accepted physics assumptions and this model are handled under the same declared-assumption framework, then checked with the same mechanized proof discipline where formalized.

\begin{table}[h]
\centering
\caption{Physics-Assumption-to-Regime-Tag Decomposition}
\label{tab:physics-regime-decomposition}
\begin{tabular}{@{}lll@{}}
\toprule
\textbf{Physics Assumption} & \textbf{Regime Decomposition} & \textbf{Justification} \\
\midrule
Landauer bound ($E \ge kT\ln 2$) & [AR], [H=cost-growth] & Declared law; yields cost lower bounds \\
Lorentz invariance & [AR] & Declared physical symmetry \\
Discrete state space & [Q\_fin], [AR] & Finite-state core; declared assumption \\
Thermodynamic constants ($\lambda, \rho$) & [AR], [H=cost-growth] & Declared conversion model \\
Planck scale ($\ell_P$) & [AR] & Declared physical constant \\
Physical-to-core encoding ($E$) & [AR] & Declared encoding map \\
\bottomrule
\end{tabular}
\end{table}

The decomposition in Table~\ref{tab:physics-regime-decomposition} makes explicit what the paper already practices: physics assumptions are declared premises in the [AR] (axiomatic regime) sense, and when they yield cost lower bounds, they inherit [H=cost-growth]. The finite-state query core [Q\_fin] captures the discrete-state assumption underlying the computational model. No physics assumption receives a special domain tag; all decompose into the existing regime taxonomy. This is a standardization claim about proof form, not a change in physical content.\claimstamp{P\ref{prop:thermo-lift},\ref{prop:landauer-constraint},\ref{prop:lorentz-discrete},\ref{prop:discrete-state-time};T\ref{thm:physical-bridge-bundle}}{AR,H=cost-growth,Q\_fin}

Coupling this to integrity yields a mechanized policy condition: when exact competence is hardness-blocked in the active regime, integrity requires abstention or explicitly weakened guarantees (Propositions~\ref{prop:integrity-resource-bound}, \ref{prop:attempted-competence-matrix}). A stronger behavioral claim such as ``abstention is globally utility-optimal for every objective'' is \emph{not} currently mechanized and remains objective-dependent.\claimstamp{P\ref{prop:integrity-resource-bound},\ref{prop:attempted-competence-matrix}}{AR}\leanmeta{\LHrng{IC}{6}{7}, \LH{IC24}, \LHrng{IC}{26}{27}, \LH{CC30}.}

\subsection{Neukart--Vinokur Thermodynamic Duality}

Neukart and Vinokur~\cite{neukart2025thermodynamic} propose a thermodynamic duality $dU = T\,dS - p\,dV + \lambda\,dC$ where $C$ is a ``complexity coordinate'' associated with computational irreversibility in spin-glass systems. Corollary~\ref{cor:neukart-vinokur} derives their duality as a specialization of our thermodynamic lift (Proposition~\ref{prop:thermo-lift}) for $\lambda > 0$.
\claimstamp{C\ref{cor:neukart-vinokur};P\ref{prop:thermo-lift}}{AR}
\leanmeta{\LH{CC50}, \LH{CC65}.}

Their spin-glass phase transitions become regime-specific instances: our [S+ETH] hard families with $k^* = n$ witness their ``complexity potential'' via explicit TAUTOLOGY reduction (Theorem~\ref{thm:dichotomy}). Physical transport via $E:\mathcal{P}\to\mathcal{D}$ extends to arbitrary substrates, not only quantum or spin systems (Theorem~\ref{thm:physical-bridge-bundle}).
\claimstamp{T\ref{thm:dichotomy},\ref{thm:physical-bridge-bundle}}{S+ETH,AR}
\leanmeta{\LH{CC16}, \LH{CT4}.}

Thus Neukart--Vinokur $\subseteq$ [AR] regime: $\lambda,\rho > 0$, explicit encoding $E:\mathcal{P}\to\mathcal{D}$ $\Rightarrow$ their thermodynamic duality holds by assumption transfer through our \coNP-complete core (Proposition~\ref{prop:physical-claim-transport}).
\claimstamp{P\ref{prop:physical-claim-transport}}{AR}
\leanmeta{\LH{CT6}.}

\subsection{Feature Selection}

In machine learning, feature selection asks which input features are relevant for prediction. This is known to be NP-hard in general~\cite{blum1997selection, guyon2003introduction}. Our results place the decision-theoretic analog in a regime map: static checking/minimization at \coNP-complete, stochastic sufficiency at PP-complete, and sequential sufficiency at PSPACE-complete.
\claimstamp{T\ref{thm:sufficiency-conp},\ref{thm:minsuff-conp},\ref{th:stochastic-sufficiency-pp},\ref{th:sequential-sufficiency-pspace}}{S,RG}\leanmeta{\LH{CC36}, \LH{CC61}, \LHrng{DC}{9}{10}.}

\subsection{Value of Information}

The value of information (VOI) framework~\cite{howard1966information,raiffa1961applied} quantifies the maximum rational payment for information. Our work addresses a different question: not the \emph{value} of information, but the \emph{complexity} of identifying which information has value.

\subsection{Model Selection}

Statistical model selection (AIC/BIC/cross-validation) provides practical heuristics for choosing among models \cite{akaike1974new,schwarz1978estimating,stone1974cross}. Our results formalize the regime-level reason heuristic selection appears: without added structural assumptions, exact optimal model selection inherits worst-case intractability, so heuristic methods implement explicit weakened-guarantee policies for unresolved structure.
\claimstamp{T\ref{thm:sufficiency-conp},\ref{thm:dichotomy},\ref{thm:tractable}}{E,S,S+ETH}\leanmeta{\LH{CC16}, \LH{CC61}, \LHrng{CC}{69}{71}.}

\subsection{Information Compression vs Certification Length}

The paper now makes an explicit two-part information accounting object (Definition~\ref{def:raw-certified-bits}): raw report bits and evidence-backed certification bits. This sharpens the information-theoretic reading of typed claims: compression of report payload is not equivalent to compression of certifiable truth conditions \cite{shannon1948mathematical,cover2006elements,slepian1973noiseless}.

From a zero-error viewpoint, this distinction is structural rather than stylistic: compressed representations can remain short while exact decodability depends on confusability structure and admissible decoding conditions \cite{shannon1956zero,korner1973graphs,lovasz1979shannon,csiszar2011information}. The same separation appears here between short reports and evidence-backed exact claims.

The rate-distortion and MDL lines make the same split in different language: there is a difference between achievable compression rate, computational procedure for obtaining/optimizing that rate, and certifiable model adequacy \cite{shannon1959coding,blahut1972computation,rissanen1978modeling,grunwald2007minimum,li2008introduction}.

Formally, Theorem~\ref{thm:exact-certified-gap-iff-admissible} and Corollary~\ref{cor:hardness-raw-only-exact} characterize the exact-report boundary: admissible exact reports require a strict certified-bit gap above raw payload, while hardness-blocked exact reports collapse to raw-only accounting. This directly links representational succinctness to certification cost under the declared contract, rather than treating ``short reports'' as automatically high-information claims.
\claimstamp{T\ref{thm:exact-certified-gap-iff-admissible};P\ref{prop:raw-certified-bit-split};C\ref{cor:hardness-raw-only-exact}}{AR}\leanmeta{\LH{CC11}, \LH{CC18}, \LH{CC21}.}

\subsection{Certifying Outputs and Proof-Carrying Claims}

Our integrity layer matches the certifying-algorithms pattern: algorithms emit candidate outputs together with certificates that can be checked quickly, separating \emph{producing} claims from \emph{verifying} claims \cite{mcconnell2010certifying,necula1997proof}. In this paper, Definition~\ref{def:solver-integrity} is exactly that soundness discipline.

At the systems level, this is the same architecture as proof-carrying code: a producer ships evidence and a consumer runs a small checker before accepting the claim \cite{necula1997proof,mcconnell2010certifying}. Our competence definition adds the regime-specific coverage/resource requirement that certifying soundness alone does not provide.

The feasibility qualifier in Definition~\ref{def:competence-regime} yields the bounded-rationality constraint used in this paper: admissible policy is constrained by computational feasibility under the declared resource model (Remark~\ref{rem:modal-should}) \cite{sep_bounded_rationality,arora2009computational}.\claimstamp{P\ref{prop:integrity-competence-separation}}{AR}
\leanmeta{\LH{IC18}.}

\paragraph{Cryptographic-verifiability perspective (scope).}
The role of cryptography in this paper is structural verifiability, not secrecy: the relevant analogy is certificate-carrying claims with lightweight verification, not confidential encoding \cite{goldwasser1989knowledge,necula1997proof}. Concretely, the typed-report discipline is modeled as a game-based reporting interface: a prover emits $(r,\pi)$ and a verifier accepts exactly when $\pi$ certifies the declared report type under the active contract. In that model, Evidence--Admissibility Equivalence gives completeness for admissible report types, while Typed Claim Admissibility plus Exact Claims Require Exact Evidence give soundness against exact overclaiming. Certified-confidence constraints then act as evidence-gated admissibility constraints on top of that verifier interface, so raw payload bits and certified bits represent distinct resources rather than stylistic notation variants.
\claimstamp{P\ref{prop:evidence-admissibility-equivalence},\ref{prop:typed-claim-admissibility},\ref{prop:exact-requires-evidence},\ref{prop:certified-confidence-gate}}{AR}\leanmeta{\LH{IC2}, \LH{IC35}, \LH{IC38}, \LH{CC75}.}

\paragraph{Three-axis integration.}
In the cited literature, these pillars are treated separately: representation-sensitive hardness, query-access lower bounds, and certifying soundness disciplines. This paper composes all three for the same decision-relevance object in one regime-typed and machine-checked framework.
\claimstamp{T\ref{thm:dichotomy};P\ref{prop:query-regime-obstruction},\ref{prop:typed-claim-admissibility}}{S+ETH,Qf,AR}\leanmeta{\LH{CC16}, \LH{CC50}, \LH{CC75}.}
